\begin{knowbox}
\textbf{Thời gian:} $60$ phút. \quad \textbf{Tổng điểm:} $10$ điểm. \quad \textbf{Số câu:} $20$.\\[2pt]\emph{Phần A -- mỗi câu $0{,}4$ điểm. \quad Phần B -- mỗi câu $0{,}5$ điểm. \quad Phần C -- mỗi câu $0{,}75$ điểm.}
\end{knowbox}
\gap
\dividerbar{Phần A -- Trả lời nhanh (4 điểm)}
\begin{probbox}{1}{\relax}
Tìm $12{,}5\%$ của $80$.
\end{probbox}
\probsol{
$80\times12{,}5:100=10$.
\dapso{$10$.}}
\gap
\begin{probbox}{2}{\relax}
Viết tỉ số phần trăm của $3$ và $8$.
\end{probbox}
\probsol{
$3:8=0{,}375=37{,}5\%$.
\dapso{$37{,}5\%$.}}
\gap
\begin{probbox}{3}{\relax}
Tính thể tích hình lập phương có cạnh $4$ cm.
\end{probbox}
\probsol{
$4\times4\times4=64$ cm$^3$.
\dapso{$64$ cm$^3$.}}
\gap
\begin{probbox}{4}{\relax}
$2{,}5$ m$^3=\ldots$ dm$^3$.
\end{probbox}
\probsol{
$1$ m$^3=1\,000$ dm$^3$, nên $2{,}5$ m$^3=2\,500$ dm$^3$.
\dapso{$2\,500$ dm$^3$.}}
\gap
\begin{probbox}{5}{\relax}
Tính diện tích toàn phần của hình hộp chữ nhật có chiều dài $6$ cm, chiều rộng $4$ cm, chiều cao $3$ cm.
\end{probbox}
\probsol{
$(6\times4+6\times3+4\times3)\times2=(24+18+12)\times2=108$ cm$^2$.
\dapso{$108$ cm$^2$.}}
\gap
\begin{probbox}{6}{\relax}
Một ô tô đi với vận tốc $48$ km/giờ trong $45$ phút. Ô tô đi được bao nhiêu ki-lô-mét?
\end{probbox}
\probsol{
$45$ phút $=0{,}75$ giờ; $48\times0{,}75=36$ km.
\dapso{$36$ km.}}
\gap
\begin{probbox}{7}{\relax}
Đi bộ $9$ km với vận tốc $4$ km/giờ thì hết bao nhiêu thời gian?
\end{probbox}
\probsol{
$9:4=2{,}25$ giờ $=2$ giờ $15$ phút.
\dapso{$2$ giờ $15$ phút.}}
\gap
\begin{probbox}{8}{\relax}
Một thuyền có vận tốc khi nước lặng là $10$ km/giờ, vận tốc dòng nước là $2$ km/giờ. Tính vận tốc của thuyền khi đi ngược dòng.
\end{probbox}
\probsol{
$10-2=8$ km/giờ.
\dapso{$8$ km/giờ.}}
\gap
\begin{probbox}{9}{\relax}
Tung một đồng xu $25$ lần, có $12$ lần xuất hiện mặt sấp. Số lần xuất hiện mặt sấp chiếm bao nhiêu phần trăm tổng số lần tung?
\end{probbox}
\probsol{
$12:25=0{,}48=48\%$.
\dapso{$48\%$.}}
\gap
\begin{probbox}{10}{\relax}
Gà và thỏ có $10$ cái đầu và $28$ cái chân. Hỏi có bao nhiêu con thỏ?
\end{probbox}
\probsol{
Giả sử tất cả là gà: $20$ chân, thiếu $8$ chân; mỗi con thỏ nhiều hơn $2$ chân: số thỏ $8:2=4$.
\dapso{$4$ con thỏ.}}
\dividerbar{Phần B -- Tính toán (3 điểm)}
\begin{probbox}{11}{\relax}
Sau khi giảm giá $20\%$, một chiếc ba lô có giá $240\,000$ đồng. Hỏi trước khi giảm giá, chiếc ba lô có giá bao nhiêu?
\end{probbox}
\probsol{
$240\,000$ đồng ứng với $100\%-20\%=80\%$ giá cũ.\par
Giá cũ: $240\,000:80\times100=300\,000$ đồng.
\dapso{$300\,000$ đồng.}}
\gap
\begin{probbox}{12}{\relax}
Một bể nước dạng hình hộp chữ nhật dài $1{,}2$ m, rộng $0{,}5$ m, cao $0{,}8$ m. Hỏi khi bể đầy thì chứa được bao nhiêu lít nước?
\end{probbox}
\probsol{
$1{,}2\times0{,}5\times0{,}8=0{,}48$ m$^3=480$ dm$^3=480$ lít.
\dapso{$480$ lít.}}
\gap
\begin{probbox}{13}{\relax}
Tính diện tích xung quanh và diện tích toàn phần của hình lập phương có cạnh $7$ cm.
\end{probbox}
\probsol{
Diện tích một mặt: $7\times7=49$ cm$^2$.\par
Diện tích xung quanh: $49\times4=196$ cm$^2$; diện tích toàn phần: $49\times6=294$ cm$^2$.
\dapso{$196$ cm$^2$; $294$ cm$^2$.}}
\gap
\begin{probbox}{14}{\relax}
Hai thành phố A và B cách nhau $210$ km. Lúc $8$ giờ, hai ô tô khởi hành từ A và từ B, đi ngược chiều nhau với vận tốc $45$ km/giờ và $60$ km/giờ. Hỏi hai xe gặp nhau lúc mấy giờ?
\end{probbox}
\probsol{
Mỗi giờ hai xe gần nhau thêm $45+60=105$ km.\par
Thời gian: $210:105=2$ giờ; hai xe gặp nhau lúc $8+2=10$ giờ.
\dapso{Lúc $10$ giờ.}}
\gap
\begin{probbox}{15}{\relax}
Một đoàn tàu dài $200$ m chạy qua một đường hầm dài $300$ m hết $25$ giây (tính từ lúc đầu tàu vào hầm đến lúc đuôi tàu ra khỏi hầm). Tính vận tốc của tàu theo km/giờ.
\end{probbox}
\probsol{
Quãng đường: $200+300=500$ m; vận tốc: $500:25=20$ m/giây.\par
$20$ m/giây $=20\times3\,600:1\,000=72$ km/giờ.
\dapso{$72$ km/giờ.}}
\gap
\begin{probbox}{16}{\relax}
Trong $35$ học sinh có $20$ bạn đi bơi, $17$ bạn học múa và $5$ bạn không tham gia hoạt động nào. Hỏi có bao nhiêu bạn tham gia cả hai hoạt động?
\end{probbox}
\probsol{
Số bạn tham gia ít nhất một hoạt động: $35-5=30$.\par
Số bạn tham gia cả hai: $20+17-30=7$.
\dapso{$7$ bạn.}}
\dividerbar{Phần C -- Giải toán (3 điểm)}
\begin{probbox}{17}{\stars{1}}
Một hình lập phương cạnh $5$ cm được sơn tất cả các mặt, rồi cắt thành các hình lập phương nhỏ cạnh $1$ cm. Có bao nhiêu hình lập phương nhỏ được sơn \pa{a} đúng $2$ mặt; \pa{b} đúng $1$ mặt; \pa{c} không được sơn mặt nào?
\end{probbox}
\probsol{
\pa{a} Mỗi cạnh trong $12$ cạnh có $5-2=3$ hình ở giữa: $12\times3=36$.\par
\pa{b} Mỗi mặt trong $6$ mặt có $3\times3=9$ hình ở phía trong: $6\times9=54$.\par
\pa{c} Phần lõi: $3\times3\times3=27$. \ (\emph{Thử lại:} $8+36+54+27=125$.)
\dapso{\pa{a} $36$;\ \pa{b} $54$;\ \pa{c} $27$.}}
\gap
\begin{probbox}{18}{\stars{1}}
Ngày thứ nhất, một cửa hàng bán $\dfrac13$ số gạo và thêm $4$ kg. Ngày thứ hai, cửa hàng bán một nửa số gạo còn lại và thêm $3$ kg. Cửa hàng còn $22$ kg gạo. Hỏi lúc đầu cửa hàng có bao nhiêu ki-lô-gam gạo?
\end{probbox}
\probsol{
Trước ngày thứ hai: $(22+3)\times2=50$ kg.\par
$50+4=54$ kg ứng với $1-\dfrac13=\dfrac23$ số gạo lúc đầu.\par
Lúc đầu: $54:\dfrac23=81$ kg.\par
\emph{Thử lại:} $81-27-4=50$; $50-25-3=22$.
\dapso{$81$ kg.}}
\gap
\begin{probbox}{19}{\stars{1}}
Một chiếc thuyền xuôi dòng từ A đến B hết $2$ giờ và ngược dòng từ B về A hết $3$ giờ. Vận tốc dòng nước là $3$ km/giờ. Tính vận tốc của thuyền khi nước lặng và độ dài quãng sông AB.
\end{probbox}
\probsol{
Vận tốc xuôi dòng hơn vận tốc ngược dòng $3\times2=6$ km/giờ.\par
Trong $2$ giờ, đi xuôi dòng được nhiều hơn đi ngược dòng $6\times2=12$ km; giờ ngược dòng thứ ba đi đúng $12$ km đó.\par
Vận tốc ngược dòng: $12$ km/giờ; vận tốc khi nước lặng: $12+3=15$ km/giờ.\par
Quãng sông AB: $12\times3=36$ km. \ (\emph{Thử lại:} $(15+3)\times2=36$.)
\dapso{$15$ km/giờ; $36$ km.}}
\gap
\begin{probbox}{20}{\stars{2}}
Một cuộc thi có $40$ câu hỏi: câu đúng được $5$ điểm, câu sai bị trừ $2$ điểm, câu bỏ trống được $0$ điểm. Thu bỏ trống $4$ câu và được $124$ điểm. Hỏi Thu trả lời đúng bao nhiêu câu?
\end{probbox}
\probsol{
Thu trả lời $40-4=36$ câu.\par
Nếu đúng cả $36$ câu: $5\times36=180$ điểm; Thu kém đi $180-124=56$ điểm.\par
Mỗi câu sai làm mất $5+2=7$ điểm: số câu sai $56:7=8$; số câu đúng $36-8=28$.
\dapso{$28$ câu.}}
